% Relatório metodológico da análise lexicométrica do corpus latouriano
% para inserção como subseção ou apêndice metodológico do capítulo 2 da tese.
% Gerado a partir das decisões registradas em docs/decisoes_metodologicas.md.

\section{Como conduzi a análise textual: passo a passo}
\label{sec:metodo-analise-figuracoes}

Apresento nesta seção a sequência de operações que realizei para sustentar empiricamente a tensão figural discutida no capítulo \ref{capitulo2}. O objetivo é tornar a análise auditável e reflexiva, em coerência com o registro etnográfico que esta tese conduz sobre as cadeias de mediação técnica. Trabalhei sobre seis textos de Bruno Latour: três livros monográficos (\emph{Laboratory Life}, escrito com Steve Woolgar em 1979 e reeditado em 1986; \emph{Science in Action}, de 1987; \emph{Pandora's Hope}, de 1999), dois artigos metateóricos (\enquote{On Actor-Network Theory: A Few Clarifications}, de 1996, e \enquote{On Recalling ANT}, de 1999) e \emph{An Inquiry into Modes of Existence}, de 2013. O ciclo de execução ocupou duas semanas de maio de 2026 e foi conduzido com auxílio de Claude Code como mediador técnico.

\subsection{Constituição do corpus}

A primeira decisão foi escolher quais obras processar. Optei pelos três livros monográficos de 1986, 1987 e 1999 por razão analítica: rastrear a trajetória figural de Latour entre o estudo etnográfico inicial do laboratório, a sistematização da teoria ator-rede e a reflexão posterior sobre realismo científico e mediação. A escolha de três obras do mesmo autor permite leitura de trajetória que análise de obra única não oferece: posso observar quais conceitos persistem entre 1986 e 1999, quais são reformulados e quais são introduzidos em momentos específicos. Trabalhei com os originais em inglês de cada obra, mantendo o francês e o português como camadas a serem mobilizadas em rodadas posteriores. A escolha pela segunda edição de \emph{Laboratory Life} (Princeton, 1986) frente à primeira (Sage, 1979) opera diferença metodológica que registro: a segunda edição suprime a palavra \enquote{Social} do subtítulo e inclui um pós-escrito em que os autores justificam essa supressão, registro de deslocamento figural que entra como objeto da própria análise.

Mais tarde, quando observei na leitura qualitativa que o vocabulário figural dos artigos metateóricos parecia operar em registro distinto do dos livros, estendi o corpus aos dois artigos de 1996 e 1999 (Etapa 2 da execução). Em sequência, quando uma versão integral do \emph{Recalling} foi obtida por canais institucionais, refiz a análise daquele artigo sobre o texto completo (Etapa 2-bis). Por fim, estendi até \emph{An Inquiry into Modes of Existence} (Etapa 3), o livro em que Latour declara estar completando a teoria ator-rede.

\subsection{Catálogo lexical e janela de leitura}

Mantenho o catálogo de termos rastreados em um único arquivo \texttt{campos\_lexicais/catalogo\_termos.yaml}, organizado em hierarquia autor $\to$ campo figurativo $\to$ variantes. A versão final que aplico aos seis textos tem dezessete campos para Latour, entre eles \texttt{inscription}, \texttt{black\_box}, \texttt{immutable\_mobile}, \texttt{actor\_network}, \texttt{translation}, \texttt{trial\_of\_strength}, \texttt{network}, \texttt{enrollment}, \texttt{spokesperson}, \texttt{construction}, \texttt{factish}, \texttt{circulating\_reference} e o campo \texttt{militar}, este último consolidando sessenta e quatro variantes do vocabulário militar-industrial (\texttt{ally}, \texttt{enemy}, \texttt{recruit}, \texttt{mobilise} e flexões britânicas e americanas, \texttt{battle}, \texttt{war}, \texttt{conquer}, \texttt{victory}, \texttt{defeat}, \texttt{troops}, \texttt{attack}, \texttt{defend}, \texttt{fortress}, \texttt{weapon}, \texttt{soldier}, \texttt{siege} e variantes morfológicas). O catálogo em YAML é mais auditável que listas planas em \texttt{.txt}: cada campo tem registro de variantes, exclusões para descartar usos triviais e nota analítica que documenta o critério de inclusão. Para \emph{An Inquiry into Modes of Existence} adicionei um catálogo paralelo de doze campos novos (\texttt{modos\_existencia}, \texttt{preposicao}, \texttt{felicidade}, \texttt{diplomacia}, \texttt{gaia}, \texttt{valores}, \texttt{trajetoria\_passe}, \texttt{instituicao}, \texttt{categorias\_dominio}, \texttt{modernos}, \texttt{alteracao}, \texttt{experiencia}), aplicado apenas ao livro de 2013 porque a leitura retroativa sobre obras de 1986 a 1999 produziria rendimento em ordem de grandeza inferior ao ruído.

A janela KWIC, isto é, o número de palavras de contexto que envolvem cada ocorrência do termo-alvo, ficou em $\pm 10$ palavras como padrão de detecção e contagem. Janelas menores truncam a unidade argumentativa da prosa teórica densa de Latour; janelas maiores que vinte palavras geram contexto longo demais para leitura rápida durante validação amostral. Para curadoria de passagens citáveis no capítulo 2 ampliei a janela para $\pm 50$ palavras em rodada específica do refinamento.

\subsection{Extração, normalização e classificação por página}

Cada PDF do corpus foi convertido em texto bruto pelo utilitário \texttt{pdftotext} com a opção \texttt{-layout}, depositado em \texttt{corpus/txt/<id>.txt}. A inspeção amostral dos três textos da Etapa 1 registrou três artefatos sistemáticos de extração que precisariam de tratamento antes da contagem: hifenização de fim de linha em \emph{Laboratory Life} (578 ocorrências do padrão \texttt{infor-\textbackslash{}nmation}), marcadores numéricos \texttt{((NN))} injetados no corpo de \emph{Science in Action} pelo conversor (276 ocorrências), e em \emph{Pandora's Hope} cerca de 2.500 caracteres não-ASCII dominados por soft hyphens U+00AD inseridos no interior de palavras (de modo que \texttt{performed} aparecia como \texttt{per\,$\langle$U+00AD$\rangle$\,formed}, e \texttt{technology} como \texttt{tech\,$\langle$U+00AD$\rangle$\,nology}). Implementei seis operações de normalização em \texttt{scripts/01b\_normalize\_text.py}, em ordem: remoção do soft hyphen, de-hifenização de fim de linha por regex que preserva hifens médios como \texttt{actor-network}, substituição dos marcadores \texttt{((NN))} por espaço, descarte de linhas com cabeçalho tipografado em letras espaçadas (\texttt{P\ A\ N\ D\ O\ R\ A'\ S\ \ H\ O\ P\ E}), aplicação de NFKC e mapeamento de aspas tipográficas e travessões longos para ASCII, e substituição de replacement chars \texttt{\textbf{?}} por espaço. Os textos crus permanecem intocados em \texttt{corpus/txt/}; a versão normalizada que entra na análise reside em \texttt{corpus/txt\_norm/}. A separação é parte do gesto metodológico: o texto cru registra o que o conversor entregou, a versão normalizada é o que o pipeline lê.

A heurística de classificação por página, em \texttt{scripts/01\_extract\_text.py}, marca cada página com uma de cinco classes: \texttt{inicio\_capitulo}, \texttt{corpo}, \texttt{notas\_fim}, \texttt{paratexto}, \texttt{qualidade\_baixa}. A classificação distingue o front matter (sumário, prefácio, agradecimentos, dedicatória, introdução) do back matter (bibliografia, índice remissivo, apêndice) e do corpo do livro, usando uma máquina de estados de três valores e regex de cabeçalhos capazes de absorver as três variações tipográficas de início de capítulo encontradas no corpus (\texttt{Chapter N}, \texttt{C\ H\ A\ P\ T\ E\ R} com letras espaçadas em \emph{Pandora's Hope}, \texttt{chapter one} por extenso). A primeira versão da heurística tinha um bug que classificava como paratexto quase todas as páginas dos três livros, dado que toda obra começa com sumário e o estado entrava em \texttt{paratexto} na primeira página sem sair. Identifiquei o problema na amostra estratificada de validação, decompus o estado em três valores (\texttt{front\_matter}, \texttt{corpo}, \texttt{back\_matter}, com \texttt{notas\_fim} como transição), e a distribuição final ficou alinhada com o conteúdo efetivo das obras: 200 páginas de corpo em \emph{Laboratory Life}, 277 em \emph{Science in Action}, 297 em \emph{Pandora's Hope}.

\subsection{Contagem, frequência e cocorrência}

A operação em torno da qual o pipeline se organiza é o KWIC, sigla de \emph{keyword in context}, gerado por \texttt{scripts/02\_kwic.py}: o script percorre o texto normalizado de cada obra, casa cada termo do catálogo via expressão regular com word boundaries, descarta ocorrências cujo contexto adjacente casa as expressões de exclusão definidas no YAML, e grava em \texttt{outputs/etapa<N>/<obra>/csv/kwic.csv} uma linha por ocorrência com contexto $\pm 10$ palavras. O script seguinte (\texttt{03\_frequencies.py}) agrega o KWIC em tabela de frequência por campo, calculando densidade por dez mil palavras a partir do total de palavras do corpo de cada obra. As visualizações por obra (barras horizontais com a ordem de densidade dos campos, histograma da posição relativa das ocorrências do campo militar ao longo do texto) ficam em \texttt{outputs/etapa<N>/<obra>/figuras/}, em PNG (300 dpi) e SVG para incorporação à tese.

A matriz de cocorrência entre campos, em \texttt{scripts/05\_cooccurrence.py}, registra para cada par de campos quantas vezes uma ocorrência de um aparece na vizinhança da outra dentro de janela configurável. Para os livros mantive a janela de 200 palavras adotada na Etapa 1; para os artigos teóricos, em que 200 palavras equivaleriam a fração elevada do total, gerei duas versões da matriz, uma com janela 200 como controle e outra com janela proporcional de 2\% do texto (97 palavras para o \emph{Recalling} integral, 157 para o \emph{Clarifications}); para \emph{An Inquiry into Modes of Existence}, em que 2\% das 194 mil palavras totais resultaria em janela exagerada para significância analítica, ajustei a janela proporcional para 0,02\%, isto é, 39 palavras, mantendo a janela 200 como controle. O grafo gerado a partir da matriz, com pesos das arestas proporcionais à frequência dos pares, está em PNG e SVG na pasta de figuras de cada obra.

A trajetória conceitual entre 1986 e 1999, consolidada por \texttt{scripts/07\_trajectory.py}, agrega as três obras da Etapa 1 em tabela única que registra densidade e contagem absoluta de cada campo em cada obra. A tabela responde a três perguntas: quais figurações aparecem nas três obras e em qual densidade relativa, quais aparecem em uma e desaparecem ou persistem nas seguintes, e como o vocabulário coautoral de \emph{Laboratory Life} se relaciona com o de Latour solo nas obras posteriores.

\subsection{Validação amostral e desambiguação}

A análise computacional só sustenta o argumento da tese se os números resistirem a inspeção qualitativa. Operei essa inspeção em duas camadas.

A primeira camada é a validação da heurística de classificação por página. Sorteei amostra estratificada de quarenta e cinco páginas, distribuídas em três páginas por estrato e por obra, e revisei cada uma para confirmar se a classe predita pelo algoritmo corresponde ao conteúdo efetivo da página. A taxa de erro confirmado ficou em 0\%, abaixo da tolerância de 20\% que havia estabelecido como gate de aprovação. Os seis casos rotulados como \enquote{parcial} (páginas em que a heurística inferiu pelo estado da vizinhança sem marcador objetivo na própria página) foram inspecionados individualmente e ficam registrados como pendência reconhecida do pipeline.

A segunda camada é a desambiguação do campo militar. A contagem bruta do campo em \emph{Pandora's Hope} (1999) registra 212 ocorrências, com predominância de \texttt{war} e \texttt{wars} (85 ocorrências conjuntas). A leitura amostral revelou que parte dessas ocorrências de \texttt{war} corresponde a usos descritivo-históricos (referências à Primeira Guerra, à Segunda Guerra, às chamadas \emph{Science Wars}) que apenas marcam contexto temporal sem mobilizar o vocabulário figural latouriano. Classifiquei manualmente cada ocorrência de \texttt{war} e \texttt{wars} em \emph{Pandora's Hope} em quatro categorias (\texttt{descritivo\_historico}, \texttt{figurativo\_latouriano}, \texttt{descritivo\_bibliografico}, \texttt{ambiguo}). A planilha resultante, \texttt{outputs/etapa1/refinamento/war\_pandora\_classificacao.csv}, é o artefato de validação e a contagem refinada do campo militar caiu 26,4\% em \emph{Pandora's Hope}, mantendo o argumento sobre o vocabulário militar-industrial e ajustando a magnitude relativa entre as três obras.

Para os artigos teóricos da Etapa 2 a desambiguação ganhou três categorias adicionais: \texttt{metalinguistico}, que marca passagens em que Latour cita o próprio vocabulário da teoria ator-rede para submetê-lo à autocrítica (a sequência \enquote{vocabulary association, translation, alliance, obligatory passage point} no \emph{Recalling}, em que o autor enumera os termos da teoria para questionar a pobreza vocabular); \texttt{descritivo\_bibliografico}, que marca ocorrências em referências bibliográficas ou em títulos de obras citadas (\enquote{La nouvelle alliance} de Prigogine e Stengers); e \texttt{conceitual\_debate}, que marca uso do termo militar para descrever polêmica entre escolas teóricas. Os gatilhos automáticos dessas três categorias estão implementados em \texttt{scripts/arquivo/15\_etapa2\_desambiguar\_militar.py}. A cobertura automática das quatro ocorrências militares dos dois artigos foi de 100\%, todas em categorias não-figurais, sustentando empiricamente a hipótese de divisão de trabalho metafórico por gênero textual: o vocabulário militar-industrial dominante nos livros monográficos recua a zero nos artigos metateóricos.

\subsection{Validação amostral semântica nos artigos: protocolo A/B/C}

Para os quatro campos que organizam o argumento têxtil-topológico nos artigos (\texttt{textil}, \texttt{topologia}, \texttt{network}, \texttt{actor\_network}) executei protocolo de validação semântica em três camadas. A camada A sorteia cinco ocorrências por campo cuja janela KWIC tem maior número de termos do mesmo campo na vizinhança imediata (heurística de densidade que privilegia candidatas a citação na tese). A camada B sorteia cinco ocorrências aleatórias por campo, com semente fixa, para representar o fundo do campo sem viés de densidade. A camada C sorteia cinco ocorrências cujas variantes são as menos frequentes do campo, sob a hipótese de que variantes raras concentram polissemia (o termo \texttt{tie} usado como verbo em vez de laço, \texttt{net} como rede de computadores em vez de tessitura).

Recebi a planilha preenchida com as oitenta e duas classificações em quinze de maio de 2026. Distribuição final: 62 ocorrências classificadas como uso figural (75,6\%), 9 como parcial (11,0\%) e 11 como não-figural (13,4\%). A taxa de figuralidade ponderada ficou em 0,815. A leitura por obra registrou contraste interno aos artigos metateóricos: o \emph{Clarifications} opera com taxa de figuralidade de 0,875, enquanto o \emph{Recalling} cai para 0,636. Dos onze \enquote{não} no corpus validado, nove são da subcategoria \texttt{metalinguistico} e oito desses nove estão no \emph{Recalling}, em passagens em que Latour cita o próprio vocabulário para suspendê-lo (\enquote{let us abandon the words actor and network}, \enquote{the third nail in the coffin}). O achado adiciona divisão de segundo nível à hipótese: dentro do gênero metateórico, o \emph{Clarifications} opera em regime expositivo-figural e o \emph{Recalling} em regime autocrítico-metalinguístico.

\subsection{Recalling integral e a passagem que faltava}

O \emph{Recalling} entrou na Etapa 2 como zip de OCR página-a-página com metade das páginas severamente truncada. A cobertura inicial era de 25,3\% do artigo, com declaração inicial errônea de cerca de 80\%. Quando obtive por canais institucionais o texto integral do artigo, refiz a análise daquele texto sobre as 4.825 palavras completas, em rodada que registrei como Etapa 2-bis. A reanálise confirmou os achados da Etapa 2 e acrescentou a passagem-chave que faltava: a sequência \enquote{vocabulary association, translation, alliance, obligatory passage point} traz uma ocorrência adicional de \texttt{alliance} classificada automaticamente como \texttt{metalinguistico}, fechando a lacuna que havia ficado registrada como ressalva metodológica. A contagem refinada figural do campo militar no \emph{Recalling} permanece em zero, agora sustentada por duas ocorrências (uma descritivo-histórica para \texttt{wars}, uma metalinguística para \texttt{alliance}) em vez de uma só.

A Etapa 2-bis também revelou um regime fluido-circulatório no \emph{Recalling} que estava ausente do corpus parcial: o par \texttt{circulating\_reference}--\texttt{topologia} aparece com nove cocorrências na janela proporcional de 97 palavras. Esse regime distingue o registro do \emph{Recalling} tanto do registro expositivo-figural do \emph{Clarifications} quanto do registro descritivo-militar dos livros monográficos, e o registro entra na seção do capítulo 2 que mobiliza a tabela \texttt{tab:textil-topologico-refinado}.

\subsection{An Inquiry into Modes of Existence: continuidade e deslocamento}

Estendi a análise a \emph{An Inquiry into Modes of Existence} (Latour, 2013, tradução de Catherine Porter) por dois motivos. O primeiro é etnográfico: a Etapa 2-bis havia documentado no \emph{Recalling} o gesto autocrítico que reduziu a contagem refinada figural do campo militar a zero, e a Etapa 3 me permite observar o que ocupa o lugar do vocabulário militar no livro que Latour declara estar completando a teoria ator-rede. O segundo é de integridade reflexiva: cobrir uma janela temporal de 1986 a 2013 amplia a base sobre a qual sustento o argumento da tensão figural no capítulo 2.

A contagem refinada do campo militar em \emph{An Inquiry into Modes of Existence} (6,33 por dez mil palavras, 123 ocorrências figurais sobre 129 brutas após desambiguação automática) coloca o livro em patamar intermediário entre os artigos metateóricos (densidade zero) e os livros monográficos solo (12,19 em \emph{Pandora's Hope}, 26,03 em \emph{Science in Action}). A inspeção amostral confirma uso predominantemente figural, em registro deslocado para o horizonte cosmopolítico-diplomático: as ocorrências do campo militar articulam-se com \texttt{modernos} (144 cocorrências em janela 200), \texttt{instituicao} (78), \texttt{topologia} (78), \texttt{diplomacia} (55) e \texttt{gaia} (45). A diplomacia como categoria operatória do livro convive em vizinhança com o vocabulário militar e o ressignifica.

A malha topológica que os artigos metateóricos de 1996 e 1999 haviam inaugurado permanece em 2013 (densidade de 26,43 por dez mil para \texttt{topologia}, 13,11 para \texttt{network}, 10,34 para \texttt{textil}). As variantes top do campo \texttt{topologia} reorganizam-se em torno de \texttt{path} (90 ocorrências) e \texttt{trajectory} (61), articulando com o campo novo \texttt{trajetoria\_passe} (433 ocorrências, par dominante com 584 cocorrências em janela 200). \emph{An Inquiry into Modes of Existence} redistribui o trabalho figurativo em três camadas: vocabulário antigo persistente, vocabulário novo específico, articulação entre as duas via cocorrência.

\subsection{O que essa sequência produz para o capítulo \ref{capitulo2}}

A análise textual sistemática produziu evidência citável para o argumento da tensão figural em três registros. O primeiro registro é tabular: a tabela de trajetória conceitual entre 1986 e 1999 (Tabela \ref{tab:trajetoria-latour}), a tabela comparativa das cinco obras da Etapa 2, a tabela militar refinada e a tabela têxtil-topológica refinada pela validação amostral A/B/C. O segundo registro é gráfico: o painel de barras horizontais que compara densidade dos dezessete campos nas três obras, o histograma de densidade do campo militar ao longo dos textos, a rede de cocorrência em \emph{Science in Action} com o campo militar conectado a \texttt{enrollment}, \texttt{translation} e \texttt{actor\_network}. O terceiro registro é qualitativo: a curadoria de passagens citáveis em janela $\pm 50$ palavras, com classificação por campo e por categoria de desambiguação, fornece o material textual com que a redação interpretativa do capítulo 2 trabalha.

A análise não substitui a leitura. A análise sistematiza a presença e a distribuição dos termos figurativos ao longo das obras de cada autor, e essa sistematização sustenta uma leitura interpretativa que já estava em curso pela trajetória de pesquisa anterior. O que a análise oferece de novo é a possibilidade de auditar a leitura: cada afirmação do capítulo 2 sobre a densidade do vocabulário militar-industrial em Latour pode ser rastreada até uma linha do KWIC ou até uma página da curadoria. A análise é, em parte, o próprio objeto da tese, porque opera como exemplo concreto da mediação técnica em pesquisa contemporânea de ciências sociais e dialoga com as questões do capítulo \ref{capitulo3} sobre cadeias sedimentadas de racionalidades técnicas.

\paragraph{Reprodutibilidade.} O pipeline está versionado em \texttt{scripts/} do repositório \texttt{analise\_figuracoes}, com mapeamento de obra para etapa em arquivo único \texttt{scripts/\_paths.py} e scripts one-shot arquivados em \texttt{scripts/arquivo/}. As decisões metodológicas datadas estão em \texttt{docs/decisoes\_metodologicas.md} e o índice cronológico de briefings consumidos em \texttt{docs/historico.md}. As sementes aleatórias estão fixadas em \texttt{seed=42} para qualquer processo estocástico (amostragem estratificada de páginas, amostragem A/B/C dos campos figurativos). Os textos crus em \texttt{corpus/txt/} e a versão normalizada em \texttt{corpus/txt\_norm/} convivem como artefatos auditáveis.
